\clearpage
\thispagestyle{empty}\SourceLeaf{249}
\pdfbookmark[0]{Dioptrische Untersuchungen}{dioptrische-title}
\begin{center}
\vspace*{0.6in}
{\Large DIOPTRISCHE UNTERSUCHUNGEN\par}
\vspace{0.6in}
VON\par
\vspace{0.5in}
{\large CARL FRIEDRICH GAUSS\par}
\vspace{0.4in}
DER KÖNIGLICHEN GESELLSCHAFT VORGELEGT MDCCCXL DEZEMBER X.\par
\vfill
\rule{0.85\textwidth}{0.3pt}\par
Abhandlungen der Königlichen Gesellschaft der Wissenschaften zu Göttingen, Band I.\par
Göttingen \textsc{mdcccxliii}.\par
\rule{0.85\textwidth}{0.3pt}
\end{center}
\clearpage
\SourceLeaf{251}
\fancyhead[L]{\small Dioptrische Untersuchungen}
\renewcommand{\Article}[1]{\par\medskip\Needspace{4\baselineskip}\pdfbookmark[1]{Art. #1}{dioptrische-art-#1}\centerline{\textbf{#1.}}\nobreak\smallskip}
\pdfbookmark[0]{Dioptrische Untersuchungen}{dioptrische-body}
\begin{center}
DIOPTRISCHE UNTERSUCHUNGEN.
\end{center}
\medskip
Die Betrachtung des Weges, den Lichtstrahlen durch Glaslinsen nehmen, wenn sie gegen deren gemeinsame Achse sehr wenig geneigt sind, und der davon abhängigen Erscheinungen liefert sehr elegante Ergebnisse. Durch die Arbeiten von \textsc{Cotes}, \textsc{Euler}, \textsc{Lagrange} und \textsc{Möbius} könnten diese Ergebnisse erschöpft scheinen, doch bleibt noch manches zu wünschen übrig. Ein wesentlicher Mangel der von diesen Mathematikern aufgestellten Sätze ist, dass sie die Dicke der Linsen vernachlässigen. Dadurch erhalten die Sätze einen ungenauen und naturwidrigen Charakter, der ihren Wert stark mindert. Ohne zu bestreiten, dass es für manche anderen dioptrischen Untersuchungen, insbesondere solche, die die sogenannte Abweichung wegen der Kugelgestalt der Linsenflächen berücksichtigen, sehr nützlich, ja notwendig ist, die Dicke der Linsen zunächst zu vernachlässigen, um einfachere und handlichere Regeln für Überschläge und erste Näherungen zu gewinnen, wird man doch gern auf eine solche Preisgabe aller Genauigkeit verzichten, wo dies ohne jeden oder ohne erheblichen Verlust an Einfachheit der Ergebnisse möglich ist. Auf einen Mangel an Präzision, der das mathematische Empfinden unangenehm berührt, stoßen wir zum Teil schon bei den ersten Begriffsbestimmungen der Dioptrik. Die Begriffe der Achse und des Brennpunkts einer Linse sind zwar genau bestimmt; das gilt aber nicht für
\SourceLeaf{252}die Brennweite, die die meisten Autoren als den Abstand des Brennpunkts der Linse von ihrem Mittelpunkt definieren. Dabei setzen sie von vornherein entweder stillschweigend voraus oder erklären ausdrücklich, dass die Dicke der Linse als unendlich klein anzusehen sei. Für wirkliche Linsen behält die Brennweite dadurch eine Unbestimmtheit von der Größenordnung der Linsendicke. Wo man es genauer nimmt, misst man jenen Abstand bald von der dem Brennpunkt nächsten Linsenfläche, bald vom sogenannten optischen Mittelpunkt der Linse, bald von dem Punkt, der genau zwischen ihrer Vorder- und Hinterfläche liegt. Von all diesen Bestimmungen unterscheidet sich wiederum der Wert, der dem Vergleich der Bildgröße eines unendlich entfernten Gegenstands mit dessen scheinbarer Größe zugrunde liegen muss. Diese letzte Bestimmung ist tatsächlich die einzig zweckmäßige.

Ich habe es daher nicht für überflüssig gehalten, diesen an sich ganz elementaren Untersuchungen einige Seiten zu widmen, vor allem um zu zeigen, dass sich bei den oben erwähnten eleganten Sätzen die Dicke der Linsen ohne Verlust an Einfachheit mitberücksichtigen lässt. Beibehalten werden soll hier nur die Beschränkung auf Strahlen, die gegen die Achse unendlich wenig geneigt sind, oder die Abweichung wegen der Kugelgestalt soll außer Betracht bleiben.

\Article{1}
Die Lage aller in dieser Untersuchung vorkommenden Punkte wird durch rechtwinklige Koordinaten $x,y,z$ bestimmt. Dabei wird vorausgesetzt, dass die Mittelpunkte der verschiedenen brechenden Flächen auf der $x$-Achse liegen, und es werden nur Lichtstrahlen betrachtet, die mit dieser Achse einen sehr kleinen Winkel bilden: Bei völlig beliebigem Ursprung werden die Koordinaten $x$ in Richtung der Lichtstrahlen als wachsend angenommen.

Zunächst betrachten wir die Wirkung einer einzigen Brechung auf den Weg eines Lichtstrahls. Das Brechungsverhältnis beim Übergang aus dem ersten Medium in das zweite sei wie $\frac1n$ zu $\frac1{n'}$ oder wie $n'$ zu $n$. Mit $M$ bezeichnen wir den Mittelpunkt der kugelförmigen Grenzfläche zwischen den beiden Medien, mit $N$ den Schnittpunkt dieser Fläche mit der ersten Koordinatenachse; zugleich sollen dieselben Buchstaben auch die diesen Punkten zugehörigen
\SourceLeaf{253}Werte von $x$ bezeichnen. So soll es im Folgenden auch bei anderen Punkten der ersten Koordinatenachse gehalten werden. Ferner sei $r=M-N$, also $r$ der Radius der Grenzfläche, positiv oder negativ, je nachdem das erste Medium auf der konvexen oder auf der konkaven Seite liegt; $P$ sei der Punkt, in dem der Lichtstrahl die Grenzfläche trifft, und $\theta$ der (spitze) Winkel zwischen $MP$ und der $x$-Achse.

Die Gerade, die ein Lichtstrahl vor der Brechung beschreibt, wird durch zwei Gleichungen bestimmt, denen wir folgende Form geben:
\[
\begin{aligned}
y&=\frac\beta n(x-N)+b\\
z&=\frac\gamma n(x-N)+c
\end{aligned}
\]
Ebenso seien die Gleichungen der Geraden, die derselbe Lichtstrahl nach der Brechung beschreibt,
\[
\begin{aligned}
y&=\frac{\beta'}{n'}(x-N)+b'\\
z&=\frac{\gamma'}{n'}(x-N)+c'
\end{aligned}
\]
Es kommt also darauf an, die Abhängigkeit der vier Größen $\beta',\gamma',b',c'$ von $\beta,\gamma,b,c$ zu bestimmen. Für den Punkt $P$ gilt
\[
x=N+r(1-\cos\theta)
\]
Da für diesen Punkt sowohl die ersten als auch die zweiten Gleichungen gelten, folgt
\[
\frac\beta n.r(1-\cos\theta)+b=\frac{\beta'}{n'}.r(1-\cos\theta)+b'
\]
und damit, weil $\beta,\beta',\theta$ als unendlich kleine Größen erster Ordnung gelten, bis auf Größen dritter Ordnung genau
\[
b'=b\quad .\quad .\quad .\quad .\quad .\quad .\quad .\quad .\quad .\quad .\quad .\qquad(1)
\]
und ebenso
\[
c'=c\quad .\quad .\quad .\quad .\quad .\quad .\quad .\quad .\quad .\quad .\quad .\qquad(1)
\]

Eine durch $M$ senkrecht zur $x$-Achse gelegte Ebene schneide den ersten, nötigenfalls verlängerten Weg des Lichtstrahls in $Q$, den zweiten in $Q'$. Da $PQ'$ mit $PQ$ und $PM$ in einer Ebene liegt, liegen $M,Q,Q'$ auf einer Geraden. Bezeichnet man mit $\lambda,\lambda'$ die Winkel dieser Geraden mit $PQ,PQ'$, so sind offenbar
\SourceLeaf{254}$MQ.\sin\lambda,MQ'.\sin\lambda'$ gleich den Produkten aus dem positiv genommenen Radius der Kugelfläche und den Sinuswerten des Einfalls- und des Brechungswinkels, also den Zahlen $n',n$ proportional. Somit gilt
\[
MQ'=\frac{n.MQ.\sin\lambda}{n'\sin\lambda'}
\]
Da nun für den Punkt $Q$
\[
\begin{aligned}
y&=b+\frac{\beta r}n\\
z&=c+\frac{\gamma r}n
\end{aligned}
\]
und für den Punkt $Q'$ dagegen
\[
\begin{aligned}
y&=b'+\frac{\beta'r}{n'}\\
z&=c'+\frac{\gamma'r}{n'}
\end{aligned}
\]
gilt und sich die beiden letzteren Koordinaten zu den beiden ersteren wie $MQ'$ zu $MQ$ verhalten, erhält man
\[
\begin{aligned}
b'+\frac{\beta'r}{n'}&=\frac{n\sin\lambda}{n'\sin\lambda'}.\left(b+\frac{\beta r}n\right)\\
c'+\frac{\gamma'r}{n'}&=\frac{n\sin\lambda}{n'\sin\lambda'}.\left(c+\frac{\gamma r}n\right)
\end{aligned}
\]
oder
\[
\begin{aligned}
\beta'&=\frac{nb+\beta r}{r}.\frac{\sin\lambda}{\sin\lambda'}-\frac{n'b'}r\\
\gamma'&=\frac{nc+\gamma r}{r}\raisebox{0.55em}{.}\frac{\sin\lambda}{\sin\lambda'}-\frac{n'c'}r
\end{aligned}
\]
Diese Ausdrücke sind streng richtig; da sich aber $\lambda,\lambda'$ vom rechten Winkel um Größen erster Ordnung und ihre Sinuswerte daher von eins um Größen zweiter Ordnung unterscheiden, gilt bis auf Größen dritter Ordnung genau
\[
\left.\begin{aligned}
\beta'&=\beta-\frac{n'-n}r.b=\beta+\frac{n'-n}{N-M}.b\\
\gamma'&=\gamma-\frac{n'-n}r.c=\gamma+\frac{n'-n}{N-M}.c
\end{aligned}\right\}\quad .\quad .\quad .\quad .\quad .\quad .\qquad(2)
\]
Diese Gleichungen (1), (2) enthalten die Lösung unserer Aufgabe.

Es verdient Beachtung, dass sich dieselben Formeln auch unmittelbar auf einen reflektierten Strahl anwenden lassen, wenn man nur $-n$ für $n'$ einsetzt. Mit diesem Verfahren lassen sich auch sämtliche folgenden Untersuchungen sehr leicht auf den Fall erweitern, dass an die Stelle der Brechungen eine oder mehrere Reflexionen treten.

\SourceLeaf{255}
\Article{2}
Um die allgemeine Aufgabe zu lösen, den Weg eines Lichtstrahls nach $(\mu+1)$ aufeinanderfolgenden Brechungen zu bestimmen, wollen wir folgende Bezeichnungen verwenden.

$N^0,N',N''\, .\, .\, .\, .\, .\, N^{(\mu)}$ sind die Punkte, in denen die brechenden Flächen die $x$-Achse schneiden.

$M^0,M',M''\, .\, .\, .\, .\, M^{(\mu)}$ sind die auf dieser Achse liegenden Mittelpunkte der brechenden Flächen

$n':n^0,\ n'':n',\ n''':n''\, .\, .\, .\, .\, .\, n^{(\mu+1)}:n^{(\mu)}$ sind die Brechungsverhältnisse beim Übergang aus dem ersten Medium (vor $N^0$) in das zweite (zwischen $N^0$ und $N'$), aus dem zweiten in das dritte und so weiter. In der Emanationstheorie sind also die Zahlen $n^0,n',n''$ und so weiter den Ausbreitungsgeschwindigkeiten des Lichts in den einzelnen Medien direkt, in der Wellentheorie umgekehrt proportional. Ist das letzte Medium dasselbe wie das erste, so gilt $n^{(\mu+1)}=n^0$.

Die Gleichungen für den Weg des Lichtstrahls vor der ersten Brechung seien
\[
\begin{aligned}
y&=\frac{\beta^0}{n^0}(x-N^0)+b^0\\
z&=\frac{\gamma^0}{n^0}(x-N^0)+c^0
\end{aligned}
\]
und die Gleichungen für seinen Weg nach der ersten Brechung
\[
\begin{aligned}
y&=\frac{\beta'}{n'}(x-N^0)+b^0\\
z&=\frac{\gamma'}{n'}(x-N^0)+c^0
\end{aligned}
\]
oder, statt auf $N^0$ auf $N'$ bezogen,
\[
\begin{aligned}
y&=\frac{\beta'}{n'}(x-N')+b'\\
z&=\frac{\gamma'}{n'}(x-N')+c'
\end{aligned}
\]
ebenso die Gleichungen für den Weg nach der zweiten Brechung
\[
\begin{aligned}
y&=\frac{\beta''}{n''}(x-N')+b'\\
z&=\frac{\gamma''}{n''}(x-N')+c'
\end{aligned}
\]
oder
\[
\begin{aligned}
y&=\frac{\beta''}{n''}(x-N'')+b''\\
z&=\frac{\gamma''}{n''}(x-N'')+c''
\end{aligned}
\]
\SourceLeaf{256}und so weiter. Bezeichnen wir also die letzten Glieder in den Folgen der $\beta,\gamma,n,N,b,c$, nämlich $\beta^{(\mu+1)},\gamma^{(\mu+1)},n^{(\mu+1)},N^{(\mu)},b^{(\mu)},c^{(\mu)}$, zur Kennzeichnung als letzte Glieder mit $\beta^*,\gamma^*,n^*,N^*,b^*,c^*$, so lauten die Gleichungen für den letzten Weg des Lichtstrahls
\[
\begin{aligned}
y&=\frac{\beta^*}{n^*}(x-N^*)+b^*\\
z&=\frac{\gamma^*}{n^*}(x-N^*)+c^*
\end{aligned}
\]
Schließlich setzen wir zur Abkürzung
\[
\left.\begin{aligned}
\frac{N'-N^0}{n'}&=t',&\frac{N''-N'}{n''}&=t'',&\frac{N'''-N''}{n'''}&=t'''\ \text{u. s. f.}\\
\frac{n'-n^0}{N^0-M^0}&=u^0,&\frac{n''-n'}{N'-M'}&=u',&\frac{n'''-n''}{N''-M''}&=u''\ \text{u. s. f.}
\end{aligned}\right\}\quad .\quad .\qquad (3)
\]
und entsprechend für die letzten Glieder dieser Folgen
\[
t^{(\mu)}=t^*,\qquad u^{(\mu)}=u^*
\]
Nach dem vorhergehenden Artikel gilt demnach
\[
\begin{aligned}
\beta'&=\beta^0+u^0b^0\\
b'&=b^0+t'\beta'\\
\beta''&=\beta'+u'b'\\
b''&=b'+t''\beta''\\
\beta'''&=\beta''+u''b''\\
b'''&=b''+t'''\beta'''
\end{aligned}
\]
und so weiter. Daraus wird ersichtlich, dass $b^*,\beta^*$ linear durch $b^0$ und $\beta^0$ bestimmt werden und dass, wenn man
\[
\left.\begin{aligned}
b^*&=gb^0+h\beta^0\\
\beta^*&=kb^0+l\beta^0
\end{aligned}\right\}\quad .\quad .\quad .\quad .\quad .\quad .\quad .\quad .\quad .\qquad (4)
\]
setzt, in der von \textsc{Euler} (\textit{Comment. Nov. Acad. Petropol.} T. IX) eingeführten Bezeichnung gilt
\[
\left.\begin{aligned}
g&=(u^0,t',u',t'',u''\ .\ .\ .\ .\ .\ t^*)\\
h&=(t',u',t'',u''\ .\ .\ .\ .\ .\ t^*)\\
k&=(u^0,t',u',t'',u''\ .\ .\ .\ .\ .\ .\ u^*)\\
l&=(t',u',t'',u''\ .\ .\ .\ .\ .\ .\ u^*)
\end{aligned}\right\}\quad .\quad .\quad .\quad .\quad .\quad .\quad .\qquad (5)
\]
\SourceLeaf{257}Diese Bezeichnung bedeutet bekanntlich: Bildet man aus einer gegebenen Folge von Größen $a,a',a'',a'''$ und so weiter nach dem folgenden Algorithmus eine andere Folge, $A,A',A'',A'''$ und so weiter,
\[
A=a,\quad A'=a'A+1,\quad A''=a''A'+A,\quad A'''=a'''A''+A'\quad\text{u. s. f.}
\]
so schreibt man
\[
A=(a),\quad A'=(a,a'),\quad A''=(a,a',a''),\quad A'''=(a,a',a'',a''')\quad\text{u. s. f.}
\]
Ferner ist unmittelbar klar, dass in den Gleichungen für die dritte Koordinate $z$ die Konstanten für den letzten Weg aus denen für den ersten genauso hergeleitet werden wie in den Gleichungen für $y$, oder dass gilt
\[
\left.\begin{aligned}
c^*&=gc^0+h\gamma^0\\
\gamma^*&=kc^0+l\gamma^0
\end{aligned}\right\}\quad .\quad .\quad .\quad .\quad .\quad .\quad .\quad .\quad .\qquad (4)
\]
Die Gleichungen (3), (5), (4) enthalten die vollständige Lösung unserer Aufgabe.

\Article{3}
\textsc{Euler} hat an der genannten Stelle die wichtigsten Beziehungen zu dem erwähnten Algorithmus hergeleitet. Hier sollen nur zwei davon in Erinnerung gerufen werden.

Erstens gilt stets
\[
(a,a',a''\ldots a^{(\lambda)})(a',a''\ldots a^{(\lambda+1)})-(a,a',a''\ldots a^{(\lambda+1)})(a',a''\ldots a^{(\lambda)})=\pm1
\]
wobei das obere oder das untere Zeichen gilt, je nachdem die Gesamtzahl aller Elemente $a,a',a''\ldots a^{(\lambda+1)}$, also die Zahl $\lambda+2$, ungerade oder gerade ist.

Zweitens darf man die Reihenfolge der Elemente umkehren; es gilt nämlich
\[
(a,a',a''\ldots a^{(\lambda)})=(a^{(\lambda)},\ldots a'',a'\ a)
\]
Wendet man den ersten dieser Sätze auf die Größen $g,h,k,l$ an, so folgt
\[
gl-hk=1
\]
Die Gleichungen (4) lassen sich daher auch folgendermaßen darstellen:
\SourceLeaf{258}
\[
\begin{aligned}
b^0&=lb^*-h\beta^*\\
\beta^0&=-kb^*+g\beta^*\\
c^0&=lc^*-h\gamma^*\\
\gamma^0&=-kc^*+g\gamma^*
\end{aligned}
\]

\Article{4}
Es sei $P$ ein gegebener Punkt auf der Geraden, die den ersten Weg des Lichtstrahls darstellt und nötigenfalls verlängert wird, und $\xi,\eta,\zeta$ seien seine Koordinaten. Es gilt also
\[
n^0\eta=\beta^0(\xi-N^0)+n^0b^0
\]
oder, wenn man für $\beta^0,b^0$ die am Schluss des vorhergehenden Artikels angegebenen Ausdrücke einsetzt,
\[
n^0\eta=(g\beta^*-kb^*)(\xi-N^0)-n^0(h\beta^*-lb^*)
\]
und folglich
\[
b^*=\frac{n^0\eta+(n^0h-g(\xi-N^0))\beta^*}{n^0l-k(\xi-N^0)}
\]
Setzt man diesen Wert in die erste Gleichung für den Weg des Lichtstrahls nach der letzten Brechung ein, also in
\[
y=\frac{\beta^*}{n^*}(x-N^*)+b^*
\]
und schreibt zur Abkürzung
\[
\begin{aligned}
N^*-\frac{n^0h-g(\xi-N^0)}{n^0l-k(\xi-N^0)}.n^*&=\xi^*\\
\frac{n^0\eta}{n^0l-k(\xi-N^0)}&=\eta^*
\end{aligned}
\]
so wird diese Gleichung zu
\[
y=\eta^*+\frac{\beta^*}{n^*}(x-\xi^*)
\]
Auf genau dieselbe Weise wird, wenn man außerdem
\[
\frac{n^0\zeta}{n^0l-k(\xi-N^0)}=\zeta^*
\]
setzt, die zweite Gleichung für den Weg des Lichtstrahls nach der letzten Brechung zu
\SourceLeaf{259}
\[
z=\zeta^*+\frac{\gamma^*}{n^*}(x-\xi^*)
\]
Der Punkt $P^*$ mit den Koordinaten $\xi^*,\eta^*,\zeta^*$ liegt also auf der Geraden, die diesen letzten Weg darstellt und nötigenfalls rückwärts verlängert wird. Da seine Koordinaten von $\beta^0,b^0,\gamma^0,c^0$ unabhängig sind, ist zugleich klar, dass dieser Punkt für \textit{alle} einfallenden Strahlen, die durch $P$ gehen, derselbe ist. Man kann den Punkt $P$ als einen Gegenstand und $P^*$ als sein Bild betrachten; der Gegenstand kann aber nur dann reell sein, wenn $P$ im ersten Medium liegt oder $\xi-N^0$ negativ ist, und ebenso ist das Bild nur dann reell, wenn $P^*$ im letzten Medium liegt oder $\xi^*-N^*$ positiv ist. In den entgegengesetzten Fällen sind Gegenstand oder Bild nur virtuell.

Die Punkte $P,P^*$ liegen mit der $x$-Achse in einer Ebene. Ihre Abstände von dieser Achse verhalten sich wie eins zu der Zahl $\frac{n^0}{n^0l-k(\xi-N^0)}$, deren positives oder negatives Vorzeichen angibt, ob die Punkte auf derselben Seite der Achse oder auf entgegengesetzten Seiten liegen. Ein System von Punkten in derselben Ebene senkrecht zur $x$-Achse kann als zusammengesetzter Gegenstand betrachtet werden. Dessen zusammengesetztes Bild liegt ebenfalls in einer Ebene senkrecht zur $x$-Achse und ist dem Gegenstand ähnlich, sodass das lineare Verhältnis entsprechender Teile durch die Zahl
\[
\frac{n^0}{n^0l-k(\xi-N^0)}=g+\frac{k}{n^*}(\xi^*-N^*)
\]
ausgedrückt wird, deren Vorzeichen die aufrechte oder umgekehrte Lage unterscheidet.

\Article{5}
Das bisher Entwickelte enthält die gesamte Theorie der Veränderungen, die der Weg der Lichtstrahlen durch Brechungen erfährt. Es lässt sich leicht auch auf den Fall ausdehnen, dass Brechungen mit einer oder mehreren Reflexionen verbunden sind; dieser Fall soll hier jedoch nicht im Einzelnen ausgeführt werden. Es ist aber nicht überflüssig, die Ergebnisse in eine andere Form zu bringen, indem man sie statt auf die erste und letzte Fläche oder auf die Punkte $N^0,N^*$ auf zwei andere Punkte $Q,Q^*$ bezieht. Es seien
\[
\begin{aligned}
y&=\frac{\beta^0}{n^0}(x-Q)+B\\
z&=\frac{\gamma^0}{n^0}(x-Q)+C
\end{aligned}
\]
\SourceLeaf{260}die Gleichungen für den ersten und
\[
\begin{aligned}
y&=\frac{\beta^*}{n^*}(x-Q^*)+B^*\\
z&=\frac{\gamma^*}{n^*}(x-Q^*)+C^*
\end{aligned}
\]
die Gleichungen für den letzten Weg des Lichtstrahls, und man setze
\[
\frac{N^0-Q}{n^0}=\theta,\qquad \frac{Q^*-N^*}{n^*}=\theta^*
\]
Wir erhalten also
\[
\begin{aligned}
b^0&=B+\theta\beta^0,&c^0&=C+\theta\gamma^0\\
B^*&=b^*+\theta^*\beta^*,&C^*&=c^*+\theta^*\gamma^*
\end{aligned}
\]
Zusammen mit den Gleichungen (4) folgt daraus leicht, dass, wenn man
\[
\begin{aligned}
G&=g+\theta^*k\\
H&=h+\theta g+\theta\theta^*k+\theta^*l\\
K&=k\\
L&=l+\theta k
\end{aligned}
\]
setzt,
\[
\begin{aligned}
B^*&=GB+H\beta^0,&C^*&=GC+H\gamma^0\\
\beta^*&=KB+L\beta^0,&\gamma^*&=KC+L\gamma^0
\end{aligned}
\]
gilt. Die Koeffizienten $G,H,K,L$, die auf diese Weise an die Stelle von $g,h,k,l$ treten, erfüllen auch die Gleichung
\[
GL-HK=1
\]

\Article{6}
Andere Punkte einzuführen, auf die man die Lage des einfallenden und des austretenden Strahls bezieht, hat den Zweck, die Abhängigkeit der letzteren Lage von der ersteren einfacher darzustellen. Dafür eignen sich besonders zwei Punktpaare, die mit $E,E^*$ und $F,F^*$ bezeichnet werden sollen. Die Werte der dabei auftretenden Größen lassen sich bequem in einer Tabelle überblicken.
\SourceLeaf{261}
\[
\renewcommand{\arraystretch}{1.8}
\begin{array}{l|l|l}
&\multicolumn{1}{c|}{\mathrm{I}}&\multicolumn{1}{c}{\mathrm{II}}\\\cline{2-3}
\theta&\dfrac{1-l}{k}&-\dfrac{l}{k}\\
\theta^*&\dfrac{1-g}{k}&-\dfrac{g}{k}\\
Q&E=N^0-\dfrac{n^0(1-l)}{k}&F=N^0+\dfrac{n^0l}{k}=E+\dfrac{n^0}{k}\\
Q^*&E^*=N^*+\dfrac{n^*(1-g)}{k}&F^*=N^*-\dfrac{n^*g}{k}=E^*-\dfrac{n^*}{k}\\
G&1&0\\
H&0&-\dfrac1k\\
K&k&k\\
L&1&0
\end{array}
\]
Bringt man also die Gleichungen für den einfallenden Strahl in die Form
\[
\begin{aligned}
y&=\frac{\beta^0}{n^0}(x-E)+B\\
z&=\frac{\gamma^0}{n^0}(x-E)+C
\end{aligned}
\]
oder in die folgende Form (wobei wir die konstanten Teile zur Unterscheidung von denen der ersten Form mit Strichen bezeichnen),
\[
\begin{aligned}
y&=\frac{\beta^0}{n^0}(x-F)+B'\\
z&=\frac{\gamma^0}{n^0}(x-F)+C'
\end{aligned}
\]
so lauten die Gleichungen für den austretenden Strahl
\[
\begin{aligned}
y&=\frac{\beta^0+kB}{n^*}.(x-E^*)+B\\
z&=\frac{\gamma^0+kC}{n^*}.(x-E^*)+C
\end{aligned}
\]
oder
\[
\begin{aligned}
y&=\frac{kB'}{n^*}.(x-F^*)-\frac{\beta^0}k\\
z&=\frac{kC'}{n^*}.(x-F^*)-\frac{\gamma^0}k
\end{aligned}
\]

\SourceLeaf{262}\Article{7}
Mithilfe der Punkte $E,E^*$ lässt sich die Abhängigkeit des letzten Weges des Lichtstrahls vom ersten einfach so ausdrücken: Der letzte Weg hat relativ zum Punkt $E^*$ dieselbe Lage, die der nur einmal gebrochene Lichtstrahl relativ zu $E$ hätte, wenn sich in $E$ eine brechende Fläche mit dem Radius $\frac{n^0-n^*}{k}$ befände, durch die der Lichtstrahl aus dem ersten Medium unmittelbar in das letzte überginge. Dies gilt für den Fall, dass das erste und das letzte Medium verschieden sind. Sind sie dagegen gleich, oder gilt $n^*=n^0$, wie bei der Brechung durch eine oder mehrere Glaslinsen, so hat der letzte Weg relativ zu $E^*$ dieselbe Lage, die er relativ zu $E$ durch die Brechung an einer in $E$ befindlichen unendlich dünnen Linse mit der Brennweite $-\frac{n^0}k$ hätte. Mit anderen Worten: Man darf den Übergang aus dem ersten Medium in das letzte durch mehrere Brechungen je nachdem, ob das erste und das letzte Medium verschieden oder gleich sind, durch eine einzige Brechung oder durch eine einzige Linse von unendlich kleiner Dicke ersetzen. Dabei gibt man im ersten Fall der brechenden Fläche den Radius $\frac{n^0-n^*}{k}$, im zweiten der Linse die Brennweite $-\frac{n^0}k$, nimmt die brechende Fläche oder die Linse in $E$ an und verschiebt in beiden Fällen die Lage des austretenden Strahls um die Entfernung des Punktes $E^*$ von $E$. Das Vorzeichen des Radius der brechenden Fläche ist dabei wie oben in Artikel 1 zu verstehen, das Vorzeichen der Brennweite so, wie weiter unten in Artikel 9 angegeben werden wird.

Wegen dieser Bedeutung scheinen die Punkte $E,E^*$ durchaus eine eigene Bezeichnung zu verdienen: Ich werde sie die \textit{Hauptpunkte} des Systems von Medien, der Linse oder des Linsensystems nennen, auf das sie sich beziehen; $E$ den ersten und $E^*$ den zweiten Hauptpunkt. Unter den Ebenen der Hauptpunkte werden die durch sie senkrecht zur $x$-Achse gelegten Ebenen verstanden werden,

\Article{8}
Für die Punkte $F,F^*$ zeigen die Formeln des Artikels 6, dass allen einfallenden Lichtstrahlen, die durch den Punkt $F$ gehen, austretende Strahlen parallel zur Achse entsprechen; einfallenden Strahlen parallel zur Achse dagegen solche austretenden Strahlen, die sich im Punkt $F^*$ schneiden; für
\SourceLeaf{263}Strahlen, die von der entgegengesetzten Seite kommen, vertauschen diese Punkte ihre Funktionen. Erweitern wir also den für einzelne Linsen bestehenden Sprachgebrauch, so können $F,F^*$ die Brennpunkte des Systems von Medien oder Linsen genannt werden, auf das sie sich beziehen, $F$ der erste und $F^*$ der zweite; die durch diese Punkte senkrecht zur $x$-Achse gelegten Ebenen können Brennpunktebenen heißen. Jene Formeln des Artikels 6 zeigen zugleich, dass allen Strahlen, die sich in irgendeinem anderen Punkt der ersten Brennpunktebene schneiden, austretende Strahlen entsprechen, die gegen die Achse geneigt, aber untereinander parallel sind. Umgekehrt entsprechen allen einfallenden Strahlen, die untereinander, aber nicht zur Achse parallel sind, solche austretenden Strahlen, die sich in einem von $F^*$ verschiedenen Punkt der zweiten Brennpunktebene schneiden.

\Article{9}
Mithilfe dieser vier Ebenen erhalten wir eine sehr einfache Konstruktion für die Lage des austretenden Strahls.

Der einfallende Strahl schneide die erste Brennpunktebene im Punkt $(1)$ und die erste Hauptebene im Punkt $(2)$; eine Parallele zu $(1)(2)$ durch $F$ treffe die erste Hauptebene in $(3)$; eine Parallele zur Achse durch $(2)$ treffe die zweite Hauptebene in $(4)$; schließlich treffe eine Parallele zur Achse durch $(3)$ die zweite Brennpunktebene in $(5)$. Dann gibt $(4)(5)$ oder $(5)(4)$ die Lage des austretenden Strahls an. Die Koordinaten haben nämlich die Werte
\[
\renewcommand{\arraystretch}{1.25}
\begin{array}{c||c|c|c}
\text{für}&x&y&z\\\hline
(1)&F&B'&C'\\
(2)&E&B&C\\
F&F&0&0\\
(3)&E&B-B'&C-C'\\
(4)&E^*&B&C\\
(5)&F^*&B-B'&C-C'
\end{array}
\]
Aus den Formeln des Artikels 6 folgt also, dass der austretende Strahl durch $(4)$ und $(5)$ geht; das Erste folgt unmittelbar, das Zweite daraus, dass
\SourceLeaf{264}
\[
\begin{aligned}
B-B'&=\frac{\beta^0}{n^0}(E-F)=-\frac{\beta^0}k\\
C-C'&=\frac{\gamma^0}{n^0}(E-F)=-\frac{\gamma^0}k
\end{aligned}
\]

Im häufigsten Fall, dass $n^*=n^0$ und damit $F^*-E^*=E-F$ gilt, wird die Konstruktion noch einfacher, weil der Punkt $(3)$ überflüssig ist; man muss nur $(1),(2),(4)$ wie zuvor bestimmen und $(4)(5)$ parallel zu $(1)E$ ziehen.

Geht die Richtung des einfallenden Strahls durch $E$, so geht die Richtung des austretenden Strahls stets durch $E^*$ und ist im Fall $n^*=n^0$ zugleich zur ersteren parallel. Bei einfachen Linsen nennt man einen solchen Strahl gewöhnlich einen Hauptstrahl.

Die Entfernungen der zweiten Brennpunktebene von der zweiten Haupt\-ebene und der ersten Hauptebene von der ersten Brennpunktebene, also die Größen $-\frac{n^*}k,-\frac{n^0}k$, könnte man die \textit{Brennweiten} des Systems von Medien nennen, wenn es nicht angemessener schiene, diese Bezeichnung auf den Fall zu beschränken, dass das letzte Medium dasselbe wie das erste ist und diese Entfernungen somit gleich sind. Um dem üblichen Sprachgebrauch zu entsprechen, sehen wir die Brennweite als positiv an, wenn dem ersten Hauptpunkt eine größere Koordinate entspricht als dem ersten Brennpunkt, sodass die Brennweite stets durch $-\frac{n^0}k=-\frac{n^*}k$ ausgedrückt wird.

\Article{10}
In den oben in Artikel 4 angegebenen Formeln für die Lage des Bildes darf man, wie leicht zu sehen ist, anstelle von $N^0,N^*$ andere Punkte einsetzen, sofern man zugleich $g,h,k,l$ durch die entsprechenden $G,H,K,L$ ersetzt. Wählen wir dazu die Hauptpunkte, so erhalten wir die folgenden Ausdrücke:
\[
\begin{aligned}
\xi^*&=E^*-\frac{n^*(E-\xi)}{n^0+k(E-\xi)}\\
\eta^*&=\frac{n^0\eta}{n^0+k(E-\xi)}\\
\zeta^*&=\frac{n^0\zeta}{n^0+k(E-\xi)}
\end{aligned}
\]
Die erste Formel kann man auch in die folgende Form bringen
\[
\frac{n^*}{\xi^*-E^*}+\frac{n^0}{E-\xi}=-k
\]
\SourceLeaf{265}Wählen wir die Brennpunkte, so erhalten wir
\[
\begin{aligned}
\xi^*&=F^*+\frac{n^0n^*}{kk(F-\xi)}\\
\eta^*&=\frac{n^0\eta}{k(F-\xi)}\\
\zeta^*&=\frac{n^0\zeta}{k(F-\xi)}
\end{aligned}
\]

Da die Formeln häufig gebraucht werden, mögen sie hier auch in der Form stehen, die sie annehmen, wenn das erste und das letzte Medium gleich sind und die Brennweite mit $\varphi$ bezeichnet wird.
\[
\begin{gathered}
\frac1{\xi^*-E^*}+\frac1{E-\xi}=\frac1\varphi\\
(\xi^*-F^*)(F-\xi)=\varphi\varphi\\
\eta^*=-\frac{\varphi\eta}{F-\xi}=-\frac{\eta(\xi^*-F^*)}{\varphi}\\
\zeta^*=-\frac{\varphi\zeta}{F-\xi}=-\frac{\zeta(\xi^*-F^*)}{\varphi}
\end{gathered}
\]

\Article{11}
Die vier Hilfspunkte $E,E^*,F,F^*$ sind im besonderen Fall $k=0$ nicht mehr anwendbar. Diese Punkte wären dann als unendlich weit von den brechenden Flächen entfernt zu betrachten. In diesem Fall kann man unmittelbar die oben zur Lösung der Hauptaufgaben angegebenen allgemeinen Formeln verwenden. Sie nehmen hier die folgende Form an.

Werden die Gleichungen für den einfallenden Strahl so geschrieben
\[
\begin{aligned}
y&=\frac{\beta^0}{n^0}(x-N^0)+b^0\\
z&=\frac{\gamma^0}{n^0}(x-N^0)+c^0
\end{aligned}
\]
so lauten die Gleichungen für den austretenden Strahl
\[
\begin{aligned}
y&=\frac{l\beta^0}{n^*}.(x-N^*)+gb^0+h\beta^0\\
z&=\frac{l\gamma^0}{n^*}.(x-N^*)+gc^0+h\gamma^0
\end{aligned}
\]
Setzt man zur Abkürzung
\[
N^*-\frac{hn^*}{l}=N^{**}
\]
oder, was wegen $gl=1$ dasselbe ist,
\SourceLeaf{266}
\[
N^*-ghn^*=N^{**}
\]
so nehmen diese Formeln eine noch einfachere Form an, nämlich
\[
\begin{aligned}
y&=\frac{l\beta^0}{n^*}.(x-N^{**})+gb^0\\
z&=\frac{l\gamma^0}{n^*}.(x-N^{**})+gc^0
\end{aligned}
\]
Für die Lage des Bildes des Punktes mit den Koordinaten $\xi,\eta,\zeta$ erhalten wir die Koordinaten
\[
\begin{aligned}
\xi^*&=N^*-ghn^*-\frac{n^*}{n^0}gg(N^0-\xi)\\
&=N^{**}-\frac{n^*}{n^0}gg(N^0-\xi)\\
\eta^*&=g\eta\\
\zeta^*&=g\zeta
\end{aligned}
\]
Daraus geht hervor, dass der Punkt auf der $x$-Achse, der nach der von uns durchgehend verwendeten Bezeichnungsweise mit $N^{**}$ zu bezeichnen ist, das Bild des Punktes $N^0$ darstellt und dass das lineare Größenverhältnis der Teile eines zusammengesetzten Bildes zum Gegenstand konstant ist, nämlich $g$ zu $1$ oder $1$ zu $l$.

\Article{12}
Der im vorigen Artikel betrachtete Fall tritt bei einem Fernrohr auf, dessen Linsen für ein weitsichtiges Auge und für deutliches Sehen unendlich weit entfernter Gegenstände eingestellt sind. Aus den obigen Formeln geht hervor, dass die Richtung des austretenden Strahls allein von der Richtung des einfallenden abhängt. Untereinander parallelen einfallenden Strahlen entsprechen also ebenfalls parallele austretende Strahlen, und die Tangente des Neigungswinkels der einfallenden Strahlen zur Achse verhält sich zur Tangente des Neigungswinkels der austretenden wie $1$ zu $l$. Die Zahl $l=\frac1g$ ist also das, was man die Vergrößerung des Fernrohrs nennt. Ihr positives oder negatives Vorzeichen gibt an, ob die Gegenstände aufrecht oder umgekehrt erscheinen. Vertauscht man die Funktionen der einfallenden und der austretenden Strahlen, indem man den Gegenständen die Okularseite zuwendet. so erscheinen sie im selben Verhältnis verkleinert, und darauf beruht das ebenso bequeme wie genaue Verfahren zur Bestimmung der Vergrößerung eines Fern\SourceLeaf{267}rohrs, das ich 1823 im 2. Band der \textit{Astronomischen Nachrichten} mitgeteilt habe.

Eine andere Methode zur Bestimmung der Vergrößerung beruht darauf, einen Gegenstand mit seinem Bild anhand des linearen Größenverhältnisses zu vergleichen. \textsc{Ramsdens} Dynameter ist nichts anderes als eine Vorrichtung, mit der sich der Durchmesser des in $N^{**}$ entstehenden Bildes der kreisrunden Begrenzung des Objektivs messen lässt. Dabei muss man sich natürlich zuerst vergewissern, dass dieses Bild tatsächlich erscheint und nicht etwa durch eine innere Blende verdeckt wird. Das Bild muss auch reell sein, wozu $ghn^*$ negativ sein muss. Bei einem \textsc{Galileis}chen Fernrohr, bei dem dieses Bild nur virtuell ist, ließe sich ein genaues Ergebnis nur mit einem mikrometrischen Mikroskop erzielen. Ein solches Mikroskop würde auch in \textit{allen} Fällen. in denen eine höhere Genauigkeit gewünscht ist, den Vorzug verdienen. Aus dem vorherigen Artikel geht außerdem hervor, dass man ebenso gut einen geeigneten, vom Objektiv entfernten Gegenstand verwenden kann, solange die Entfernung nur nicht so groß wird, dass das Bild nicht mehr reell oder mit dem Mikroskop nicht mehr erreichbar ist. Schließlich sei noch bemerkt, dass $N^{**}$ der Punkt ist, der in der Theorie der Fernrohre als \textit{Ort des Auges} bezeichnet wird.

\Article{13}
Um die allgemeinen Regeln des Artikels 2 auf eine einfache Glaslinse anzuwenden, bezeichnen wir das Brechungsverhältnis beim Übergang von Luft in Glas mit $n:1$, die Radien der ersten und zweiten Fläche mit $(n-1)f$ und $(n-1)f'$ und die Dicke der Linse mit $ne$. Anstelle der dortigen Bezeichnungen haben wir also
\[
\begin{array}{ccc}
n^0&\mathit{hier}&1\\
n'&.\,.\,.\,.&n\\
n''&.\,.\,.\,.&1\\
t'&.\,.\,.\,.&e\\
u^0&.\,.\,.\,.&-\dfrac1f\\\noalign{\vskip 4pt}
u'&.\,.\,.\,.&-\dfrac1{f'}
\end{array}
\]
und folglich
\SourceLeaf{268}
\[
\begin{aligned}
g&=1+u^0t'=\frac{f-e}{f}\\
h&=t'=e\\
k&=u^0+u'+t'u^0u'=-\frac{f+f'-e}{ff'}\\
l&=1+u't'=\frac{f'-e}{f'}
\end{aligned}
\]
Für die Brennweite $\varphi$ erhalten wir also nach Artikel 9
\[
\varphi=\frac{ff'}{f+f'-e}
\]
für die beiden hier mit $E,E'$ zu bezeichnenden Hauptpunkte nach Artikel 6
\[
\begin{aligned}
E&=N^0+\frac{ef}{f+f'-e}=N^0+\frac{e\varphi}{f'}\\
E'&=N'-\frac{ef'}{f+f'-e}=N'-\frac{e\varphi}{f}
\end{aligned}
\]
und für die beiden Brennpunkte $F,F'$
\[
\begin{aligned}
F&=E-\varphi=N^0-\frac{f(f'-e)}{f+f'-e}\\
F'&=E'+\varphi=N'+\frac{f'(f-e)}{f+f'-e}
\end{aligned}
\]
Für den Schnittpunkt der Geraden, die ein Hauptstrahl im Inneren der Linse beschreibt, mit der Achse findet man leicht den folgenden Ausdruck, wobei die Gerade nötigenfalls vorwärts oder rückwärts verlängert wird
\[
x=N^0+\frac{nef}{f+f'}=N'-\frac{nef'}{f+f'}
\]
Diesen Punkt, der also von der Neigung des Hauptstrahls unabhängig ist, nennen manche Autoren den optischen Mittelpunkt der Linse. Diese Auszeichnung dürfte der Punkt, der sonst gar keine bemerkenswerten Eigenschaften hat, kaum verdienen. Die Bezeichnung scheint hier und da zu dem Irrtum verleitet zu haben, die einfachen Beziehungen zwischen Bild und Gegenstand, die bei einer unendlich dünnen Linse gelten, ließen sich allein durch den Bezug auf jenen Mittelpunkt auf eine Linse endlicher Dicke übertragen. Wie oben gezeigt, gilt diese Übertragung jedoch nur dann, wenn der Gegenstand auf den ersten und das Bild auf den zweiten Hauptpunkt bezogen wird. Bei einem System aus mehreren Linsen, also schon bei einem achromatischen Doppelobjektiv, kann ohnehin von einem Mittelpunkt in jenem Sinne gar keine Rede sein. Will man die Bezeichnung beibehalten, so würde ich es für angemessener halten, sie dem Punkt zu geben, der genau zwischen den beiden
\SourceLeaf{269}Hauptpunkten (und damit auch zwischen den beiden Brennpunkten) liegt und mit jenem Punkt nur dann zusammenfällt, wenn die Linse gleichseitig ist. Dieser Punkt hat die praktisch nützliche Eigenschaft, dass man ihn durch Umwenden der Linse leicht und genau bestimmen kann. Denn offenbar muss gerade dieser Punkt beim Umwenden wieder seinen früheren Platz einnehmen, wenn die Lage des Bildes eines festen Gegenstands unverändert bleiben soll.

Hier sei noch bemerkt, dass der Abstand der beiden Hauptpunkte
\[
E'-E=ne-\frac{e(f+f')}{f+f'-e}=(n-1)e-\frac{ee}{f+f'-e}
\]
beträgt und sich daher, insofern $e$ gegenüber $f+f'-e$ gewöhnlich sehr klein ist, von $(n-1)e$ oder von der mit $\frac{n-1}{n}$ multiplizierten Dicke der Linse kaum merklich unterscheidet.

\Article{14}
Für ein System von Linsen auf einer gemeinsamen Achse lassen sich die allgemeinen Formeln des Artikels 2, mit denen aus dem Weg des einfallenden Lichtstrahls der Weg des austretenden bestimmt wird, durch bequemere ersetzen. Dazu führt man anstelle der Radien der einzelnen brechenden Flächen und ihrer gegenseitigen Abstände die Brennweiten der einzelnen Linsen und die Abstände ihrer zweiten Hauptpunkte von den ersten Hauptpunkten der folgenden Linsen ein. Die neuen Formeln sind denen des Artikels 2 sehr ähnlich, enthalten aber nur halb so viele Elemente. Da sie sich sehr leicht aus dem Vorhergehenden ableiten lassen, genügt es, sie hier in unmittelbar verwendbarer Form anzugeben.

Wir bezeichnen die Brennweiten der einzelnen aufeinanderfolgenden Linsen mit $\varphi^0,\varphi',\varphi''$ usw. Für ihre Hauptpunkte weichen wir hier von der bisherigen Bezeichnungsweise ab. Wir bezeichnen die ersten mit $E^0,E',E''$ usw. und die zweiten mit $I^0,I',I''$ usw. Zur Abkürzung schreiben wir
\[
\begin{aligned}
-\frac1{\varphi^0}&=u^0,&-\frac1{\varphi'}&=u',&-\frac1{\varphi''}&=u''\quad\text{usw.}\\
E'-I^0&=t',&E''-I'&=t'',&E'''-I''&=t'''\quad\text{usw.}
\end{aligned}
\]
Die letzten Glieder in diesen Reihen können als solche durch ein Sternchen gekennzeichnet werden.

\SourceLeaf{270}Bringt man nun die Gleichungen für den einfallenden Strahl in die Form
\[
\begin{aligned}
y&=\beta^0(x-E^0)+b^0\\
z&=\gamma^0(x-E^0)+c^0
\end{aligned}
\]
und die für den austretenden Strahl in die folgende Form
\[
\begin{aligned}
y&=\beta^*(x-I^*)+b^*\\
z&=\gamma^*(x-I^*)+c^*
\end{aligned}
\]
so erhält man, wenn die vier Größen $g,h,k,l$ durch Formeln bestimmt werden, die mit den in Artikel 2 als (5) bezeichneten völlig übereinstimmen,
\[
\begin{aligned}
b^*&=gb^0+h\beta^0,&c^*&=gc^0+h\gamma^0\\
\beta^*&=kb^0+l\beta^0,&\gamma^*&=kc^0+l\gamma^0
\end{aligned}
\]
Für die beiden Hauptpunkte des Linsensystems als Ganzes erhält man
\[
\begin{aligned}
\text{für den ersten}\quad x&=E^0-\frac{1-l}{k}\\
\text{für den zweiten}\quad x&=I^*+\frac{1-g}{k}
\end{aligned}
\]
Ferner erhält man für die beiden Brennpunkte des Linsensystems
\[
\begin{aligned}
\text{für den ersten}\quad x&=E^0+\frac{l}{k}\\
\text{für den zweiten}\quad x&=I^*-\frac{g}{k}
\end{aligned}
\]
\[
\text{die Brennweite selbst ist}\quad=-\frac1k
\]

Die Formeln für den Fall, dass das System nur aus zwei Linsen besteht, verdienen es, eigens angegeben zu werden. Man erhält nämlich
\[
\begin{aligned}
g&=\frac{\varphi^0-t'}{\varphi^0}\\
h&=t'\\
k&=-\frac{\varphi^0+\varphi'-t'}{\varphi^0\varphi'}\\
l&=\frac{\varphi'-t'}{\varphi'}
\end{aligned}
\]
Die Werte von $x$ für die beiden Hauptpunkte sind
\SourceLeaf{271}
\[
E^0+\frac{t'\varphi^0}{\varphi^0+\varphi'-t'}\quad\text{und}\quad I'-\frac{t'\varphi'}{\varphi^0+\varphi'-t'}
\]
und die Brennweite ist
\[
=\frac{\varphi^0\varphi'}{\varphi^0+\varphi'-t'}
\]
Man sieht, dass diese Formeln denen ganz entsprechen, die in Artikel 13 zur Bestimmung der Hauptpunkte und der Brennweite einer einfachen Linse angegeben wurden. An die Stelle der dortigen $f^0,f',e$ treten hier die Größen $\varphi^0,\varphi',t'$.

Der Abstand der beiden Hauptpunkte beträgt im Fall zweier Linsen
\[
\begin{aligned}
&=I'-E^0-\frac{t'(\varphi^0+\varphi')}{\varphi^0+\varphi'-t'}\\
&=I^0-E^0+I'-E'-\frac{t't'}{\varphi^0+\varphi'-t'}
\end{aligned}
\]
Ist $t'$ sehr klein, wie es bei achromatischen Doppellinsen der üblichen Bauart immer der Fall ist, so wird das letzte Glied unbedeutend. Der Abstand der beiden Hauptpunkte einer solchen Doppellinse als Ganzes ist daher fast genau gleich der Summe der beiden Werte, die dieser Abstand in den einzelnen Linsen hat.

Es ist übrigens unmittelbar klar, dass sich sämtliche in diesem Artikel angegebenen Formeln ohne jede Änderung auf den Fall übertragen lassen, dass statt einfacher Linsen Teilsysteme von Linsen zu einem Gesamtsystem vereinigt werden sollen.

\Article{15}
Die optischen Erscheinungen bei einer einfachen Linse wie auch bei einem System mehrerer Linsen auf einer gemeinsamen Achse hängen, wie wir gezeigt haben, von drei Elementen ab. Diese werden durch das Brechungsverhältnis (oder durch die Brechungsverhältnisse, wenn sie für die verschiedenen Linsen verschieden sind) sowie durch die Lagen und Radien der brechenden Flächen bestimmt. Da diese Größen gewöhnlich nicht unmittelbar bekannt sind, bleibt noch einiges über die Methode zu sagen, mit der man umgekehrt jene drei Elemente aus beobachteten Erscheinungen ableiten kann. Wir bezeichnen die verschiedenen dabei in Frage kommenden Punkte der Achse wie folgt:
\SourceLeaf{272}$\xi$ ist ein Gegenstand und $\xi'$ sein Bild. $F$ und $F'$ sind der erste und der zweite Brennpunkt, $E$ und $E'$ der erste und der zweite Hauptpunkt, und schließlich ist $D$ ein Punkt, der mit der Linse (oder dem Linsensystem) fest verbunden ist. Dieselben Buchstaben bezeichnen wie immer die Koordinaten dieser Punkte in jedem Versuch. Wir setzen ferner die Brennweite $=f$ und schreiben für die Abstände des Punktes $D$ von den Brennpunkten $D-F=p$ und $F'-D=q$. Die drei Größen $f,p,q$ kann man als die Elemente der Linse betrachten. Zu ihrer Bestimmung sind daher immer drei Versuche erforderlich. Man muss bei drei verschiedenen Lagen des Gegenstands und seines Bildes relativ zur Linse deren Abstände vom Punkt $D$ messen. Diese Aufgabe wollen wir zunächst ganz allgemein lösen.

Die Werte von $D-\xi$ und $\xi'-D$ seien in einem Versuch $a,b$, in einem zweiten $a',b'$ und in einem dritten $a'',b''$. Die allgemeine Gleichung
\[
(F-\xi)(\xi'-F')=ff
\]
liefert uns demnach
\[
\begin{aligned}
(a-p)(b-q)&=ff\\
(a'-p)(b'-q)&=ff\\
(a''-p)(b''-q)&=ff
\end{aligned}
\]
woraus man durch Elimination leicht erhält
\[
\begin{aligned}
p&=a-\frac{(a'-a)(a''-a)(b'-b'')}{R}\\
q&=b-\frac{(b-b')(b-b'')(a''-a')}{R}\\
ff&=\frac{(a'-a)(a''-a)(a''-a')(b-b')(b-b'')(b'-b'')}{RR}
\end{aligned}
\]
wobei zur Abkürzung
\[
(a''-a)(b-b')-(a'-a)(b-b'')=R
\]
geschrieben wird. Man kann $R$ auch in die folgende Form bringen
\[
\begin{aligned}
R&=(a''-a')(b-b')-(a'-a)(b'-b'')\\
&=(a''-a')(b-b'')-(a''-a)(b'-b'')
\end{aligned}
\]
und $p$ und $q$ in die folgenden Formen
\SourceLeaf{273}
\[
\begin{aligned}
p&=a'-\frac{(a'-a)(a''-a')(b-b'')}{R}\\
&=a''-\frac{(a''-a)(a''-a')(b-b')}{R}\\
q&=b'-\frac{(b-b')(b'-b'')(a''-a)}{R}\\
&=b''-\frac{(b-b'')(b'-b'')(a'-a)}{R}
\end{aligned}
\]

\Article{16}
Der im vorigen Artikel angegebenen allgemeinen Lösung sind noch einige Bemerkungen hinzuzufügen.

I. Es wird vorausgesetzt, dass der Gegenstand in allen drei Versuchen auf derselben Seite der Linse liegt. Hält man es für zweckmäßig, die Linse in einem der Versuche umzukehren, so kann man sich diesen Versuch so vorstellen, als wäre das Bild der Gegenstand und der Gegenstand das Bild. Dadurch wird dieser Fall auf den vorigen zurückgeführt.

II. Für sich genommen lässt die Formel für $ff$ offen, ob $f$ positiv oder negativ anzusetzen ist. Dies entscheidet sich aber bereits daran, ob das Bild aufrecht oder umgekehrt ist, da $\xi'-F'$ und $f$ im ersten Fall entgegengesetzte und im zweiten gleiche Vorzeichen haben müssen. Dabei ist zu beachten, dass die praktische Anwendung bei aller Allgemeingültigkeit der analytischen Lösung auf reelle Bilder (also bei einzelnen Linsen auf positive Brennweiten) beschränkt bleibt, sofern man keine besonderen Hilfsmittel zur Bestimmung der Lage virtueller Bilder heranzieht.

III. Da die Versuche immer nur einen begrenzten Grad an Genauigkeit zulassen, hängt die Zuverlässigkeit der Ergebnisse sehr davon ab, welche Kombinationen man wählt. Im Allgemeinen kann als Regel gelten, dass sich durch drei Versuche, von denen zwei unter Umständen, die sich nur wenig voneinander unterscheiden, ausgeführt werden, jedenfalls nicht alle drei Elemente genau bestimmen lassen.

\Article{17}
Bei einer einfachen Linse ebenso wie bei einer Linse aus zwei oder mehr sehr nahe beieinanderliegenden Linsen (wie bei achromatischen Objektiven üblicher Bauart) haben die beiden Haupt\SourceLeaf{274}punkte einen geringen Abstand voneinander. Könnte man diesen Abstand $E'-E=\lambda$ als bekannte Größe betrachten, so würden \textit{zwei} Versuche genügen, da die Gleichung
\[
p+q=2f+\lambda
\]
an die Stelle des dritten Versuchs tritt. Verbindet man mit ihr die beiden anderen Gleichungen
\[
\begin{aligned}
(a-p)(b-q)&=ff\\
(a'-p)(b'-q)&=ff
\end{aligned}
\]
so erhält man nach der Elimination von $p$ und $q$ zur Bestimmung von $f$ die Gleichung
\[
\begin{aligned}
\frac{(a'+b'-a-b)^2}{(a'-a)(b-b')}.ff+2(a+b+a'+b'-2\lambda)f\\
{}-(a+b'-\lambda)(a'+b-\lambda)=0
\end{aligned}
\]
Diese quadratische Gleichung wird linear, wenn $a'+b'-a-b=0$ gilt, das heißt, wenn die beiden Versuche so angeordnet sind, dass der Abstand des Bildes vom Gegenstand in beiden derselbe bleibt, während die Linse zwei verschiedene Positionen einnimmt. Dieser Abstand sei $=c$, also $a=c-b$ und $a'=c-b'$. Daraus ergibt sich
\[
4(c-\lambda)f=(c-\lambda+b'-b)(c-\lambda-b'+b)
\]
oder
\[
f=\tfrac14(c-\lambda)-\frac{(b'-b)^2}{4(c-\lambda)}
\]
Für jeden vorgegebenen Wert von $c$ muss $F-\xi$ der Gleichung
\[
F-\xi+\frac{ff}{F-\xi}=F-\xi+\xi'-F'=c-2f-\lambda
\]
genügen. deren zwei Wurzeln
\[
\begin{aligned}
F-\xi&=\tfrac12(c-2f-\lambda)+\tfrac12\BareRoot{(c-4f-\lambda)(c-\lambda)}\\
F-\xi&=\tfrac12(c-2f-\lambda)-\tfrac12\BareRoot{(c-4f-\lambda)(c-\lambda)}
\end{aligned}
\]
reell und verschieden sind, wenn $c$ größer als $4f+\lambda$ ist. Dann gibt es für einen festen Gegenstand $\xi$ immer zwei verschiedene Positionen der Linse, bei denen das Bild mit dem Punkt $\xi+c$ zusammenfällt. Das Produkt dieser beiden Werte von $F-\xi$, also $(a-p)(a'-p)$, ist $=ff$, woraus zugleich hervorgeht,
\SourceLeaf{275}dass $a'-p=b-q$ und $b'-q=a-p$ gelten, folglich
\[
\begin{aligned}
p&=\tfrac12(2f+c+\lambda-b-b')\\
q&=\tfrac12(2f-c+\lambda+b+b')\\
E&=D+\tfrac12(b+b'-c)-\tfrac12\lambda\\
E'&=D+\tfrac12(b+b'-c)+\tfrac12\lambda
\end{aligned}
\]

\Article{18}
In der Stellung der Linse, in der $F-\xi=f$ gilt, ist $\xi'-\xi=4f+\lambda$, oder das Bild hat den kleinsten Abstand vom Gegenstand. Es entfernt sich von ihm, sobald man die Linse aus dieser Stellung nach der einen oder anderen Seite verschiebt, anfangs aber offenbar sehr langsam. Wird für $c$ ein Wert gewählt, der $4f+\lambda$ nur wenig übersteigt, so lassen die Versuche zur Bestimmung der beiden erforderlichen Linsenstellungen oder der Werte von $b$ und $b'$ folglich nur eine vergleichsweise geringe Genauigkeit zu. Diese Unsicherheit wirkt sich in vollem Umfang auf die Bestimmung von $E$ und $E'$ aus. Zu \textit{diesem} Zweck ist das Verfahren unter solchen Umständen daher kaum brauchbar. Anders verhält es sich, wenn es nur um die Bestimmung der Brennweite geht. Die Genauigkeit nimmt durch jenen Umstand überhaupt nicht ab, da in den Ausdruck für $f$ nur das Quadrat von $b'-b$ eingeht. Das Verfahren ist in diesem Fall zudem bequemer durchzuführen, da außer dem Abstand $c$ nur die Verschiebung der Linse um $b'-b$ gemessen werden muss. Die absoluten Werte von $b$ und $b'$ sind also nicht nötig.

\Article{19}
Vernachlässigt man $\lambda$ ganz und setzt also
\[
f=\tfrac14c-\frac{(b'-b)^2}{4c}
\]
so stimmt das beschriebene Verfahren mit dem überein, das \textsc{Bessel} im 17. Band der Astronomischen Nachrichten vorgeschlagen und zur Bestimmung der Brennweite des Objektivs des Königsberger Heliometers angewandt hat. Die exakte Formel zeigt, dass man bei Vernachlässigung von $\lambda$ die Brennweite um
\SourceLeaf{276}
\[
\tfrac14\lambda+\frac{\lambda(b'-b)^2}{4c(c-\lambda)}
\]
zu groß erhält. Das zweite Glied kann unter den genannten Umständen als unmerklich betrachtet werden. Für ein Ergebnis, das der mit dem Verfahren an sich möglichen Genauigkeit entspricht, muss man $\lambda$ daher unbedingt berücksichtigen. Diese Größe genau zu bestimmen ist allerdings mit einigen Schwierigkeiten verbunden. Bei einer einfachen Linse genügt es, aus ihrer gemessenen Dicke und dem ungefähr bekannten Brechungsverhältnis den in Artikel 13 angegebenen Näherungswert für $\lambda$ zu berechnen. Auch bei einer achromatischen Doppellinse kann man gegebenenfalls den in Artikel 14 genannten Näherungswert verwenden, sofern man die Dicke jedes einzelnen Bestandteils genau bestimmen kann. Um wenigstens eine ungefähre Vorstellung davon zu bekommen, welchen Einfluss die Vernachlässigung von $\lambda$ haben kann, betrachten wir als Beispiel ein Objektiv, bei dem die Kronglaslinse 7 Linien und die Flintglaslinse 3 Linien dick ist. Für die erste nehmen wir ein Brechungsverhältnis von 1,528 an und für die zweite eines von 1,618. Damit beträgt der Abstand der beiden Hauptpunkte näherungsweise
\[
\begin{array}{lr}
\text{für die Kronglaslinse}&2{,}42\\
\text{für die Flintglaslinse}&1{,}15\\
\text{und für die zusammengesetzte Linse}&3{,}57\ \text{Linien}
\end{array}
\]
sodass die Brennweite um 0,89 Linien zu groß ausfällt. Bei einem Objektiv mit 8 Fuß Brennweite und den vorausgesetzten Abmessungen würde der Fehler also etwa $\frac1{1300}$ des Ganzen betragen.

\Article{20}
Kann man das im vorigen Artikel angegebene Verfahren zur Bestimmung von $\lambda$ nicht anwenden oder will man sich damit nicht begnügen, so scheint der folgende Weg am besten geeignet, um die Größe durch unmittelbare Versuche zu bestimmen.

Man bestimme die Lage des Bildes eines sehr weit entfernten Gegenstands (so gut wie möglich auf der Achse der Linse) relativ zum Punkt $D$. Soweit man die Entfernung des Gegenstands als unendlich groß betrachten kann, fällt das Bild
\SourceLeaf{277}in $F'$, und der gemessene Abstand $F'-D$ ergibt somit unmittelbar $q$. Man wiederhole den Versuch mit umgekehrter Linse. Das Bild fällt dann in $F$, und sein Abstand von $D$ ergibt den Wert von $p$. Für den dritten Versuch bringe man den Gegenstand (auf der Seite von $F$) möglichst nahe an die Linse. Man bestimme den Abstand des Bildes von diesem Gegenstand $=\xi'-\xi$ und zugleich den Abstand $D-\xi=a''$ und setze $\xi'-D=\xi'-\xi-a''=b''$. Man erhält also
\[
(p-a'')(q-b'')=ff
\]
oder
\[
\lambda=p+q-2\BareRoot{(p-a'')(q-b'')}
\]
Konnte man die Messungen in allen drei Versuchen mit größter Genauigkeit durchführen, so sind allein dadurch alle drei Elemente $p,q,f$ hinreichend genau bestimmt, und es sind keine weiteren Versuche nötig. Möchte man aber $f$ noch genauer erhalten, so sind diese Versuche nur als Vorbereitung auf das Verfahren des Artikels 18 zu betrachten. Diese Vorbereitung liefert den Wert von $\lambda$. Um klarer zu sehen, wovon die Genauigkeit der so erhaltenen Bestimmung von $\lambda$ hauptsächlich abhängt, bringen wir die obige Formel für $\lambda$ in die folgende Form
\[
\lambda=a''+b''+\frac{(p-a''-\BareRoot{(p-a'')(q-b'')})^2}{p-a''}
\]
und beachten, dass $p-a''-\BareRoot{(p-a'')(q-b'')}$ den Abstand des Gegenstands im dritten Versuch vom ersten Hauptpunkt darstellt, $p-a''$ hingegen den Abstand dieses Gegenstands vom ersten Brennpunkt. Daraus geht hervor, dass das letzte Glied der Formel für $\lambda$ unter den gegebenen Umständen sehr klein wird und sein berechneter Wert durch kleine Ungenauigkeiten in den Werten von $p,q,a'',b''$ nur wenig beeinflusst wird. Die Genauigkeit der Bestimmung von $\lambda$ hängt daher hauptsächlich von der Genauigkeit ab, mit der $\xi'-\xi=a''+b''$ gemessen wird.

\Article{21}
Zum Verfahren des vorigen Artikels verdienen hier noch einige Bemerkungen Platz.

I. Für den dritten Versuch, bei dem das Bild nur virtuell ist, reichen die sonst brauchbaren Mittel nicht aus. Das folgende Verfahren verei\SourceLeaf{278}nt jedoch einfache Handhabung und Genauigkeit. Auf einer ebenen Fläche zeichnet man einen Kreis, der ebenso groß oder etwas größer ist als der vorspringende Rand der Linsenfassung, und kennzeichnet den Mittelpunkt durch zwei feine, sich kreuzende Linien. Man legt die Linse mit ihrer Fassung so auf die Fläche, dass die Fassung konzentrisch zum Kreis liegt. Dann stellt man ein zusammengesetztes Mikroskop mit Fadenkreuz auf einem festen Stativ senkrecht darüber und verschiebt es in seiner Hülse, bis das Bild des Linienkreuzes genau mit dem Fadenkreuz zusammenfällt. Schließlich nimmt man die Linse weg und nähert das Mikroskop durch Verschieben in der Hülse der Ebene, bis das Bild des Linienkreuzes wieder mit dem Fadenkreuz des Mikroskops zusammenfällt. Diese letzte Verschiebung lässt sich mit irgendeinem geeigneten Mittel leicht genau messen. Sie entspricht dem Abstand des Gegenstands (des Linienkreuzes) von seinem durch die Glaslinie erzeugten Bild $=\xi'-\xi$. Den Punkt auf der Linsenachse, der in der den vorspringenden Rand der Fassung berührenden Ebene liegt, kann man als festen Punkt $D$ selbst wählen. In diesem Fall gilt $a''=0$, $b''=\xi'-\xi$. Hat man einen anderen Punkt $D$ gewählt, kann man ihn mit jenem Punkt durch naheliegende Mittel vergleichen, um $a''$ zu finden.

II. Ist die Entfernung des für den ersten und zweiten Versuch benutzten Gegenstands zwar groß, aber nicht groß genug, um sein Bild als genau mit dem Brennpunkt zusammenfallend anzusehen, so ist eine Korrektur nötig. Man erhält sie, indem man das Quadrat der Brennweite durch die Entfernung des Gegenstands teilt. Diese Korrektur zieht man von den Abständen des Bildes vom Punkt $D$ ab, um die Größen $q$ und $p$ genau zu erhalten. Dafür braucht man offenbar nur eine grobe Näherung für die Brennweite und die Entfernung, sofern letztere sehr groß ist. Ebenso leicht lässt sich diese Korrektur jedoch durch direkte, streng exakte Formeln bestimmen. Für den ersten Versuch sei $a$ der Wert von $D-\xi$ und $b$ der Wert von $\xi'-D$. Für den zweiten Versuch hingegen, bei dem die Linse in umgekehrter Stellung verwendet wird, bezeichne man den Abstand $D-\xi$ mit $b'$ und $\xi'-D$ mit $a'$. Dieses Vorgehen läuft auf dasselbe hinaus wie das in Artikel 16, I, angegebene und bietet uns den Vorteil, dass die Gleichungen für die drei Versuche
\SourceLeaf{279}
\[
\begin{aligned}
(a-p)(b-q)&=ff\\
(b'-q)(a'-p)&=ff\\
(a''-p)(b''-q)&=ff
\end{aligned}
\]
mit denen übereinstimmen, von denen wir in Artikel 15 ausgingen. Damit bleiben auch die durch Elimination daraus gewonnenen Formeln gültig. Ist man beim zweiten Versuch so vorgegangen, dass das Bild im Raum dieselbe Lage hat wie im ersten (was leicht zu erreichen ist, obwohl sich das Ergebnis bei der vorausgesetzten großen Entfernung des Gegenstands nicht merklich ändert, wenn man es damit nicht so genau nimmt), so gilt $a+b=b'+a'$. Bezeichnen wir diese Größe mit $c$, so vereinfachen sich dadurch die Formeln aus Artikel 15 noch etwas. Aus der zweiten Formel für $p$ und der ersten für $q$ ergibt sich nämlich
\[
\begin{aligned}
p&=a'-\frac{(a'-a'')(b-b'')}{c-a''-b''}\\
q&=b-\frac{(a'-a'')(b-b'')}{c-a''-b''}
\end{aligned}
\]

III. Auch wenn man das Verfahren aus Artikel 20 nicht zur vollständigen Bestimmung der Elemente verwenden, sondern die genaueste Bestimmung der Brennweite der Methode aus Artikel 17 vorbehalten möchte, bleibt das Verfahren aus Artikel 20 doch am besten geeignet, die Lage der beiden Hauptpunkte festzulegen. Es ergibt sich nämlich
\[
\begin{aligned}
E&=D+\tfrac12(q-p)-\tfrac12\lambda=D+\tfrac12(b-a')-\tfrac12\lambda\\
E'&=D+\tfrac12(q-p)+\tfrac12\lambda=D+\tfrac12(b-a')+\tfrac12\lambda
\end{aligned}
\]

\Article{22}
Einer einfachen Linse und, allgemeiner gesprochen, auch einem Linsensystem kann man nur dann eine bestimmte Brennweite und bestimmte Lagen der Haupt- und Brennpunkte zuordnen, wenn es um Lichtstrahlen mit bestimmter Brechbarkeit geht. Bei Strahlen anderer Brechbarkeit liegen diese Punkte anders, und die Brennweite hat einen anderen Wert. Nicht homogenes Licht von Gegenständen wird daher beim Durchgang durch Gläser farblich zerstreut. Diese Farbzerstreuung lässt sich durch die Kombination von zwei oder
\SourceLeaf{280}mehr Linsen aus verschiedenen Glasarten aufheben. Ein vollkommen achromatisches Objektiv setzt jedoch voraus, dass sich parallele Strahlen unabhängig von ihrer Farbe in einem einzigen Punkt vereinigen. Das gilt sowohl für achsenparallele als auch für schräg zur Achse einfallende Strahlen. Anders gesagt, die verschiedenfarbigen Bilder eines ausgedehnten, als unendlich weit entfernt betrachteten Gegenstands müssen in derselben Ebene liegen und gleich groß sein. Die erste Bedingung beruht darauf, dass der zweite Brennpunkt für verschiedenfarbige Strahlen derselbe ist, die zweite darauf, dass die Brennweiten gleich sind. Da die Brennweite der Abstand des zweiten Brennpunkts vom zweiten Hauptpunkt ist, kann man die beiden Bedingungen auch so ausdrücken, dass beide Punkte für rote und violette Strahlen dieselben sein müssen. Ist nur die erste Bedingung erfüllt, so liefern die zur Achse geneigten Strahlen kein reines Bild. Ein \textit{sehr} kleiner Unterschied zwischen den Brennweiten für verschiedenfarbige Strahlen darf jedoch stets als völlig unschädlich gelten.

In der Theorie der achromatischen Objektive berücksichtigt man gewöhnlich nur die erste Bedingung. Bei der üblichen Bauart dieser Objektive berühren sich die beiden Linsen jedoch oder liegen äußerst dicht beieinander. Die unterschiedliche Brechbarkeit der Lichtstrahlen beeinflusst daher die Lage der Hauptpunkte so wenig, dass die zweite Bedingung von selbst erfüllt ist. Auch wenn sie nicht genau erfüllt ist, ist die Annäherung so gut, dass daraus keine merkliche Unvollkommenheit entstehen kann. Hält man es für der Mühe wert, so kann man außerdem die Dicken der Linsen so berechnen, dass der zweite Hauptpunkt für Strahlen unterschiedlicher Brechbarkeit genau derselbe ist.

Anders verhält es sich, wenn die konvexe Kronglaslinse von der konkaven Flintglaslinse durch einen beträchtlichen Abstand getrennt ist. Es lässt sich leicht zeigen, dass bei Werten für diesen Abstand und die Brennweiten der einzelnen Linsen, für die der zweite Brennpunkt des Linsensystems bei verschiedenfarbigen Strahlen derselbe bleibt, die Brennweite des Systems für violette Strahlen zwangsläufig größer ist als für rote. Der Unterschied hat dieselbe Größenordnung wie der bei einfachen Linsen, aber die umgekehrte Richtung. Das gilt auch dann noch, wenn man, wie bei den sogenannten dialytischen Fernrohren, die zweite Linse durch
\SourceLeaf{281}eine Kombination aus einer Flintglaslinse und einer Kronglaslinse ersetzt, die sich berühren oder in sehr geringem Abstand voneinander liegen. Auf diese Weise bleibt es stets unmöglich, von einem ausgedehnten Gegenstand ein vollkommen farbreines Bild zu erzeugen. Soll das violette Bild in derselben Entfernung vom Linsensystem liegen wie das rote, so ist es zwangsläufig größer als dieses.

Man darf daraus keineswegs schließen, dass \textit{Fernrohre} der zuletzt beschriebenen Bauart hinsichtlich der Achromasie unvollkommener bleiben müssen als Fernrohre mit achromatischen Objektiven herkömmlicher Bauart, die ein vollkommen farbreines Bild erzeugen. Vielmehr lässt sich gerade das Gegenteil behaupten. Bei einer gut berechneten Anordnung der Okulare können jene Fernrohre dem Auge das farbreinere Bild liefern.

Tatsächlich kann ein vom Objektiv erzeugtes vollkommen farbreines Bild, ob reell oder virtuell, dem \textit{Auge} wegen der von den Okularlinsen verursachten Farbzerstreuung nicht vollkommen rein \textit{erscheinen}. Durch eine besondere Anordnung der Okulare verhindert man zwar den sogenannten farbigen Rand, kann damit aber die Längenabweichung nicht beseitigen. Diese wird noch dadurch vergrößert, dass das menschliche Auge selbst nicht achromatisch ist. Man erreicht nur, dass die letzten Bilder, das rote und das violette, scheinbar gleich groß sind. Man erreicht damit weder, dass sie in gleicher Entfernung liegen, noch, dass sie \textit{zugleich} scharf erscheinen.

Die ungleiche Größe der ersten Bilder, des roten und des violetten, ist bei dialytischen Objektiven unvermeidlich. Durch eine geeignete Gestaltung der Okulare lässt sie sich aber sehr gut ausgleichen, sodass der sichtbare farbige Rand ebenso gut beseitigt wird wie bei Fernrohren üblicher Bauart. Die eben erwähnte zweite Unvollkommenheit bleibt auch hier bestehen, solange das erste rote und das erste violette Bild in gleicher Entfernung vom Objektiv liegen.

Es ist also klar, dass zur Erzeugung eines vollkommen farbreinen Bildes \textit{im Auge} das erste Bild eine gewisse Längenabweichung haben \textit{muss}. Diese Abweichung hängt von den Eigenschaften der Okulare und von der fehlenden Achromasie des menschlichen Auges ab. Theoretisch lässt sich auch ein Objektiv üblicher Bauart so berechnen, dass eine vorgegebene Längenabweichung entsteht. Abgesehen von der Schwierigkeit, bei der Herstellung die ganze
\SourceLeaf{282}Genauigkeit zu erreichen, die zur Erzeugung so sehr kleiner Unterschiede nötig ist, würde diese Längenabweichung aber immer nur für ein bestimmtes Okular passen. Bei der dialytischen Bauart lässt sich dagegen die Doppellinse, die den zweiten Teil des Objektivs bildet, gegen den ersten Teil verschieben. So erhält man die für jedes Okular erforderliche Längenabweichung, während das Okular so gestaltet werden kann, dass der farbige Rand in der nötigen Weise beseitigt wird. Im Übrigen muss ich mich hier auf diesen kurzen Hinweis beschränken und eine ausführlichere Behandlung dieses interessanten Gegenstands einer anderen Gelegenheit vorbehalten.

